
\documentclass[a4paper,11pt]{article}

\usepackage{amsmath, amssymb, amsfonts}
\usepackage{geometry}
\geometry{margin=1in}

\title{Christoffel Symbols and Example on an Elliptic Paraboloid}
\author{}
\date{}

\begin{document}
	
	\maketitle
	
	\section*{Introduction}
	
	We proceed step by step: first the logic of Christoffel symbols, then a concrete computation on the elliptic paraboloid
	\[
	(t,v,t^2+v^2).
	\]
	
	\section{What Christoffel Symbols Represent}
	
	We consider
	\[
	\Gamma^k_{ij},
	\]
	the Christoffel symbols of the second kind.
	
	\subsection*{Intuition}
	
	They encode how the \textbf{basis vectors change from point to point} on a curved manifold.
	
	In Euclidean space with Cartesian coordinates, basis vectors are constant, so
	\[
	\Gamma^k_{ij} = 0.
	\]
	
	On a curved surface (or curvilinear coordinates), the basis vectors vary, hence
	\[
	\Gamma^k_{ij} \neq 0.
	\]
	
	\section{Formal Definition}
	
	Given a metric $g_{ij}$, the Christoffel symbols are defined by
	\[
	\Gamma^k_{ij}
	=
	\frac{1}{2} g^{k\ell}
	\left(
	\partial_i g_{j\ell}
	+
	\partial_j g_{i\ell}
	-
	\partial_\ell g_{ij}
	\right),
	\]
	where:
	\begin{itemize}
		\item $g_{ij}$ is the metric tensor,
		\item $g^{ij}$ is the inverse metric,
		\item $\partial_i = \frac{\partial}{\partial x^i}$.
	\end{itemize}
	
	\section{Geometric Viewpoint}
	
	They describe covariant differentiation:
	\[
	\nabla_{\partial_i} \partial_j = \Gamma^k_{ij} \partial_k.
	\]
	
	They:
	\begin{itemize}
		\item define a connection,
		\item describe parallel transport,
		\item appear in geodesic equations:
	\end{itemize}
	
	\[
	\frac{d^2 x^k}{ds^2}
	+
	\Gamma^k_{ij}
	\frac{dx^i}{ds}
	\frac{dx^j}{ds}
	= 0.
	\]
	
	\section{Elliptic Paraboloid}
	
	Consider the parametrization
	\[
	\mathbf{X}(t,v) = (t, v, t^2 + v^2).
	\]
	
	\subsection{Step 1: Tangent Vectors}
	
	\[
	\mathbf{X}_t = (1, 0, 2t),
	\qquad
	\mathbf{X}_v = (0, 1, 2v).
	\]
	
	\subsection{Step 2: Metric Tensor}
	
	\[
	g_{ij} = \langle \mathbf{X}_i, \mathbf{X}_j \rangle.
	\]
	
	Compute:
	\[
	g_{tt} = 1 + 4t^2, \quad
	g_{tv} = 4tv, \quad
	g_{vv} = 1 + 4v^2.
	\]
	
	Thus:
	\[
	g =
	\begin{pmatrix}
		1 + 4t^2 & 4tv \\
		4tv & 1 + 4v^2
	\end{pmatrix}.
	\]
	
	\subsection{Step 3: Inverse Metric}
	
	\[
	\det g = (1+4t^2)(1+4v^2) - (4tv)^2 = 1 + 4t^2 + 4v^2.
	\]
	
	Let
	\[
	D = 1 + 4t^2 + 4v^2.
	\]
	
	Then:
	\[
	g^{-1} =
	\frac{1}{D}
	\begin{pmatrix}
		1 + 4v^2 & -4tv \\
		-4tv & 1 + 4t^2
	\end{pmatrix}.
	\]
	
	\subsection{Step 4: Derivatives of the Metric}
	
	\[
	\partial_t g_{tt} = 8t, \quad \partial_v g_{tt} = 0,
	\]
	\[
	\partial_t g_{tv} = 4v, \quad \partial_v g_{tv} = 4t,
	\]
	\[
	\partial_t g_{vv} = 0, \quad \partial_v g_{vv} = 8v.
	\]
	
	\section{Computation of Christoffel Symbols}
	
	\subsection{Examples}
	
	\subsubsection*{$\Gamma^t_{tt}$}
	
	\[
	\Gamma^t_{tt}
	=
	\frac{4t}{D}.
	\]
	
	\subsubsection*{$\Gamma^t_{tv} = \Gamma^t_{vt}$}
	
	\[
	\Gamma^t_{tv}
	=
	\frac{4v}{D}.
	\]
	
	\subsubsection*{$\Gamma^t_{vv}$}
	
	\[
	\Gamma^t_{vv}
	=
	-\frac{4t}{D}.
	\]
	
	\subsection{Second Coordinate}
	
	\[
	\Gamma^v_{tt}
	=
	-\frac{4v}{D},
	\]
	\[
	\Gamma^v_{tv}
	=
	\frac{4t}{D},
	\]
	\[
	\Gamma^v_{vv}
	=
	\frac{4v}{D}.
	\]
	
	\section{Final Result}
	
	With
	\[
	D = 1 + 4t^2 + 4v^2,
	\]
	we obtain:
	\[
	\begin{aligned}
		\Gamma^t_{tt} &= \frac{4t}{D}, &
		\Gamma^t_{tv} &= \frac{4v}{D}, &
		\Gamma^t_{vv} &= -\frac{4t}{D} \\
		\\
		\Gamma^v_{tt} &= -\frac{4v}{D}, &
		\Gamma^v_{tv} &= \frac{4t}{D}, &
		\Gamma^v_{vv} &= \frac{4v}{D}.
	\end{aligned}
	\]
	
	\section{Interpretation}
	
	\begin{itemize}
		\item The Christoffel symbols are nonzero because the surface is curved.
		\item They vanish at $(0,0)$, reflecting local flatness.
		\item Their structure reflects the symmetry of the paraboloid.
	\end{itemize}
	
	\section{Big-Picture Intuition}
	
	Christoffel symbols can be understood as:
	\begin{itemize}
		\item correction terms to ordinary derivatives,
		\item objects encoding curvature indirectly,
		\item arising from non-straight coordinate systems.
	\end{itemize}
	
\end{document}